% =========================================================================
% 生成式人工智能使用声明（按 2026 年竞赛规定置于参考文献之前）
% 队员提交前务必逐条核实工具名称、版本与使用范围，并保留使用记录备查。
% =========================================================================

\section*{生成式人工智能使用声明}
\addcontentsline{toc}{section}{生成式人工智能使用声明}

本文在准备过程中使用了生成式人工智能工具作为辅助。使用的具体情况如下：

\begin{enumerate}
  \item \textbf{工具名称与版本}：Kimi（WorkBuddy 智能体，月之暗面）与
        Codex（OpenAI），访问时间为 2026 年 9 月 10 日至 13 日。
  \item \textbf{使用范围与方式}：
    \begin{itemize}
      \item 辅助编写与调试 Python 求解代码（线性规划装配、回放框架、
            结果导出脚本），包括定位运行错误与重构重复代码；
      \item 辅助生成与修改论文插图（MATLAB 绘图脚本的样式调整与重绘）；
      \item 辅助检查论文文字与公式的一致性（如公式符号与代码实现的对照、
            表格数字与结果文件的核对），并对部分段落作语言润色。
    \end{itemize}
  \item \textbf{未使用情形}：模型总体设计、算法选择与参数协议、全部数值
        实验的执行与结果分析、结论的得出，均由作者独立完成，\textbf{不是}
        由人工智能生成。
  \item \textbf{核实责任}：人工智能输出的全部代码、数字与文字均经作者
        逐条核实、独立复算（含不调用调度代码的账本复算与 SHA-256 复现
        清单）后采信。作者对本文全部内容（包括使用人工智能辅助完成的
        部分）承担全部责任。
\end{enumerate}
